\subsection{R-29: prescribed product moments extend the common stable range to $2m-2$ rows}

\paragraph{Statement R-29.}
Let $n>m\geq3$ and $4\leq r\leq2m-2$. Put $a=n-m$,
and fix a partition $\nu\vdash ns$, $s\geq1$, with exactly
$r$ positive parts and
\[
 \nu_1>as.
\]
Define $\eta=\nu-\nu_r(1,\ldots,1)$ and use the enhanced
cones and ambient representation $W$ of G-22 in rank $r$.
These two cones have a common stable point. They also have a
common stable point $w$ with
\begin{equation}\label{R29:no_scaling}
 g w=cw,\quad g\in\mathrm{SL}_r,\quad c\in\mathbb C^*
       \quad\Longrightarrow\quad c=1.
\end{equation}
Consequently their invariant coordinate rings have the
dimensions and exact single-label extraction in G-22, and
both multiplicities occur at every integer stretch
$j\geq\beta^2$ for any common ambient generator-degree bound
$\beta$. In particular the stable-point and scalar-period
hypotheses of the higher-row Hilbert comparison extend from
$r\leq m$ to $r\leq2m-2$ under the stated strict conditions.

\paragraph{Background and normalization.}
Use the tensor-induced Hermitian metrics on symmetric powers.
For a nonzero degree-$m$ form $f$ in $k$ variables, write
$M(f)$ for its normalized Hermitian moment matrix before
subtracting the scalar part. Thus its expectation on a
Hermitian infinitesimal matrix $B$ is
$\langle Bf,f\rangle/\langle f,f\rangle
 =\operatorname{tr}(BM(f))$, and $\operatorname{tr}M(f)=m$.
Multiplying $f$ by a nonzero scalar does not change $M(f)$.
The traceless part is zero exactly when the usual
$\mathrm{SL}_k$ moment map is zero. The equivalence between
closed orbit and an orbit meeting moment zero is the cited
Kempf--Ness criterion: Losev (2006), Theorem~4, printed p.~1
(G-B14). Hilbert--Mumford and stable-locus openness are the
same cited inputs as G-22 and R-28. We do not assume that
the moment matrix of a product is the sum of its factors'
moment matrices.

\begin{proof}
We first prove a prescribed-moment lemma and then construct
the common tensor from two orthogonal variable blocks.

\paragraph{1. A product of $m$ linear forms can realize any positive rational spectrum in at most $m-1$ variables.}
Let $2\leq k\leq m-1$, and let
$b_1,\ldots,b_k\in\mathbb Q_{>0}$ have sum $m$.
There is a unit degree-$m$ form $P$, a product of $m$
distinct linear forms in $\mathbb C^k$, such that
\begin{equation}\label{R29:prescribed_moment}
 M(P)=\operatorname{diag}(b_1,\ldots,b_k),
\end{equation}
and $[P]$ has finite projective stabilizer.

Permute coordinates first so the $b_i$ are weakly decreasing.
Choose a positive integer $d$ with
$\lambda=(db_1,\ldots,db_k)$ integral, and put
$\theta=\lambda-\lambda_k(1,\ldots,1)$.
For a primitive $m$th root $\zeta$ take the unit Fourier
vectors
\[
 u_j=\sum_{i=0}^{k-1}\sqrt{b_{i+1}/m}\,\zeta^{ij}x_i,
       \qquad 0\leq j<m.
\]
As in G-22, their rank-one moment sum is
$\operatorname{diag}(b_1,\ldots,b_k)$, and every $k$ of
them are independent by the Vandermonde determinant.
Since $m\geq k+1$, they contain a projective frame.
The ordered tensor
\[
 \xi=(u_0\otimes\cdots\otimes u_{m-1})^{\otimes d}
                   \otimes v_\theta
\]
has moment zero for $\mathrm{SL}_k$: its full diagonal moment
is $\lambda-\theta=\lambda_k(1,\ldots,1)$.
Its projective stabilizer is finite, since fixing every
ordered tensor line fixes the projective frame. The cited
Kempf--Ness criterion makes its orbit closed, so $\xi$ is
stable.

Put $P_0=\prod_j u_j$. For every one-parameter subgroup,
the smallest weight of $P_0$ is the sum of the smallest
weights of the $u_j$, because their initial terms have
nonzero product in a polynomial ring. The same statement
holds for the outer power. Thus
$P_0^{\otimes d}\otimes v_\theta$ and $\xi$ have identical
smallest weight for every subgroup. Hilbert--Mumford, in
exactly the form used in G-22, proves that the former
tensor is stable as well.

Apply the reverse direction of Kempf--Ness to this closed
orbit. Some $g\in\mathrm{SL}_k$ sends
$P_0^{\otimes d}\otimes v_\theta$ to moment zero.
The line of $g v_\theta$ is a point of the partial flag
variety parametrizing the extremal orbit. A unitary matrix
can take this flag to the standard flag: choose an
orthonormal basis adapted to its nested subspaces. Hence
there is a unitary $U$ such that the normalized vector
$Ugv_\theta$ is a scalar of absolute value one times
$v_\theta$. A scalar adjustment makes $U$ special unitary
if desired, without changing these lines or moments.

Let $P$ be the unit normalization of $UgP_0$.
Unitary equivariance preserves moment zero, and the pure
outer tensor power contributes $dM(P)$. The auxiliary
factor contributes $-\theta$. Therefore
\[
 dM(P)-\operatorname{diag}(\theta)
           =cI_k.
\]
Taking traces gives $c=\lambda_k$, and proves
\eqref{R29:prescribed_moment}. Undo the initial coordinate
permutation to obtain the prescribed order of the $b_i$.
The factor lines of $P$ still contain a projective frame,
since $Ug$ is invertible. A transformation preserving
$[P]$ permutes its $m$ distinct linear factors, and the
kernel of that finite permutation action fixes the frame
pointwise. It is scalar. Thus the projective stabilizer
of $P$ is finite, completing the lemma.

\paragraph{2. Two blocks give the desired full-rank moment matrix.}
Use coordinates $x_0,\ldots,x_{r-1}$ and put
\[
 \mu_0=\nu_1/s-a,\qquad
 \mu_i=\nu_{i+1}/s\quad(1\leq i<r).
\]
These are positive rational numbers with sum $m$.
Split $V=V_1\oplus V_2$ orthogonally, with
$k_1=\lfloor r/2\rfloor$, $k_2=\lceil r/2\rceil$,
using the first $k_1$ coordinates for $V_1$.
Both block dimensions lie between two and $m-1$.
Define
\[
 w_b=\frac1m\sum_{x_i\in V_b}\mu_i\quad(b=1,2).
\]
Then $w_1,w_2>0$ and $w_1+w_2=1$. Within block $b$,
the numbers $\mu_i/w_b$ have sum $m$. By the lemma choose
unit products $P\in\operatorname{Sym}^mV_1$ and
$Q\in\operatorname{Sym}^mV_2$, each with finite projective
stabilizer, and with those prescribed moment matrices.
Set
\begin{equation}\label{R29:two_block_q}
 q=\sqrt{w_1}P+\sqrt{w_2}Q.
\end{equation}
The two pure-block spaces are orthogonal, so $q$ is unit.
Its block-diagonal moment is the corresponding weighted
sum. An off-diagonal infinitesimal operator has bidegree
$(m-1,1)$ or $(1,m-1)$, orthogonal to both pure blocks.
It follows that
\begin{equation}\label{R29:full_moment}
 M(q)=\operatorname{diag}(\mu_0,\ldots,\mu_{r-1}).
\end{equation}
This is a direct moment computation for a sum in orthogonal
subspaces, not an additive moment assertion for a polynomial
product.

The projective stabilizer of $q$ is finite. First, the
constant-coefficient directional derivatives of each block
product are independent. If $D_vP=0$, restrict the equation
to the zero hyperplane of one factor, away from the others;
it gives that this factor vanishes on $v$. Doing this for
every factor and using their span shows $v=0$.
The argument for $Q$ is identical.

Now suppose $Bq=cq$. Because $m\geq3$, the two mixed
bidegrees $(m-1,1)$ and $(1,m-1)$ are different from
one another and from the pure bidegrees. The off-diagonal
blocks of $B$ consequently vanish, by the derivative
independence just proved. The remaining equations are
$B_1P=cP$ and $B_2Q=cQ$. The finite projective stabilizers
make both $B_i$ scalar, and their degree-$m$ actions force
the same scalar $c/m$ in both blocks. Thus $B$ is scalar.
In characteristic zero this proves finiteness of the
projective algebraic stabilizer of $q$.

Write the two products in \eqref{R29:two_block_q} as
$\prod_{j=1}^m L_j$ and $\prod_{j=1}^m M_j$, absorbing
the two nonzero scalar coefficients into individual factors.
The diagonal-plus-cycle pencil with diagonal $L_j$ and
successive cycle entries $M_j$ has permanent equal to
their sum. Changing the sign of one cycle entry when
necessary gives the same sum as a determinant. Hence
$q$ belongs to both size-$m$ images.

\paragraph{3. Padding and the flag preserve stability by weights.}
Take $\ell=x_0$ and the unit extremal vector $v_\eta$.
In the ordered representation consider
\[
 \widetilde\xi=
  \bigl(x_0^{\otimes a}\otimes q\bigr)^{\otimes s}
                    \otimes v_\eta.
\]
Equations \eqref{R29:full_moment} and
$sa e_0+s\mu-\eta=\nu_r(1,\ldots,1)$ show that its
$\mathrm{SL}_r$ moment is zero. Its projective stabilizer
is finite because it preserves the line of its $q$ factor,
whose projective stabilizer is finite. Thus
$\widetilde\xi$ is stable.

Put $F=x_0^a q$ and $w=F^{\otimes s}\otimes v_\eta$.
For every one-parameter subgroup, the initial term of
$x_0^a q$ is the nonzero product of the respective initial
terms. The minimum-weight identity of G-22 therefore
applies to $\widetilde\xi,w$, and proves stability of $w$.
No moment formula for the padded polynomial product has
been used. Block-diagonal padding proves that $F$ belongs
to both required degree-$n$ cones. This gives the common
stable point throughout the stated row range.

\paragraph{4. The enlarged common family also removes the scalar period.}
Use the irreducible affine parameter space of all linear
forms $L_1,\ldots,L_m,M_1,\ldots,M_m,\ell$, and set
\[
 q=\prod_jL_j+\prod_jM_j,\qquad
 w=(\ell^a q)^{\otimes s}\otimes v_\eta,
\]
with the auxiliary vector fixed. Every point is still
represented by the same two-matching cycle construction.
The stable enhanced locus is nonempty by the preceding
steps, and the finite-projective-stabilizer locus for $q$
is nonempty by \eqref{R29:two_block_q}. Both are open.
The irreducible locus for $q$ is also nonempty: specialize
to $x_1^m+x_2^{m-1}x_3$, Eisenstein at $x_3$ as a
polynomial in $x_1$. It remains irreducible after adjoining
the unused variables, since its quotient polynomial ring
remains an integral domain. This gives a witness inside
the same parameter space even when $r>m$.

These three nonempty open sets intersect. Fix an
irreducible $q$ with finite projective automorphism group
at a stable enhanced point, and vary the padding line
inside the nonempty stable open of that fiber. Avoiding
the finitely many proper fixed loci of the nonidentity
automorphisms makes the padded form projectively free,
exactly as in R-28. Its nonlinear irreducible factor is
distinguished from the padding line even when $a=1$.
The central action has exponent $ns-|\eta|=r\nu_r$,
so the argument of R-28 proves \eqref{R29:no_scaling}.
The same closed-orbit separation and numerical-semigroup
argument gives gcd one and all stretches $j\geq\beta^2$.

Finally the enhanced-cone dimensions and graded invariant
extraction in G-22 only require a common stable point,
the rank-$r$ Cartan coordinate ring, and the stated grading.
Their proofs are independent of $r\leq m$ once stability
is supplied. They therefore apply unchanged in this
larger range.
\end{proof}

\paragraph{Scope.}
The prescribed-moment lemma is an existence construction
using Kempf--Ness normalization; it does not assert a
polynomial bit-complexity algorithm for the required
change of coordinates. The two-block method reaches
$r\leq2m-2$, and makes no claim beyond that range.
For $r>m$ and integer $n>m$, the strict rectangular
condition $n<rm/(r-1)$ is empty. The new row range is
therefore a nonrectangular strict-direction extension;
the first-row equality face is not settled by this result.
